\documentclass[12pt,a4paper,openright,twoside]{report}
\usepackage[italian]{babel}
\usepackage[latin1]{inputenc}
\usepackage{graphicx}
\usepackage{amssymb}
\usepackage{amsmath}
\usepackage{latexsym}
\usepackage{amsthm}
\oddsidemargin=30pt \evensidemargin=20pt
\hyphenation{sil-la-ba-zio-ne pa-ren-te-si}
\linespread{1.3}
\usepackage{tgtermes}
\begin{document}
	\begin{titlepage}
		\begin{center}
			
			% Nome ufficiale dell'Ateneo: NON modificare.
			{{\Large{\textsc{Alma Mater Studiorum $\cdot$ Universit\`a di
							Bologna}}}} \rule[0.1cm]{15.8cm}{0.1mm}
			\rule[0.5cm]{15.8cm}{0.6mm}
			
			% ------------------------------------------------------------------
			% DIPARTIMENTO / CORSO DI LAUREA
			% ------------------------------------------------------------------
			% Selezionare il corso di Laurea di riferimento, cancellando le opzioni non necessarie.
			%
			{\small{\bf TODO Dipartimento di Informatica --- Scienza e Ingegneria\\
					Corso di Laurea Triennale in Ingegneria e Scienze informatiche}}
		\end{center}
		
		\vspace{15mm}
		\begin{center}
			% ------------------------------------------------------------------
			% TITOLO DELLA TESI
			% ------------------------------------------------------------------
			% Sostituire "TITOLO DELLA TESI" con il titolo completo della propria tesi.
			% Suggerimenti:
			% - Usare MAIUSCOLO per il titolo principale, se richiesto dalle linee guida.
			% - Evitare abbreviazioni non standard.
			{\LARGE{\bf TODO}}\\
			\vspace{3mm}
			
			\vspace{19mm}
			% ------------------------------------------------------------------
			% SETTORE SCIENTIFICO DI RIFERIMENTO (OPZIONALE)
			% ------------------------------------------------------------------
			% Questa riga è opzionale. Puoi:
			% - Sostituire "Settore Scientifico di Riferimento" con il nome dell'insegnamento
			%   o ambito disciplinare da cui nasce la tesi (es. "Machine Learning", "Calcolo Numerico").
			% - Oppure, se non richiesto dal tuo Corso di Laurea, puoi COMMENTARE
			%   l'intera riga anteponendo un % all'inizio.
			{\large{\bf TODO}}
		\end{center}
		
		\vspace{30mm}
		\par
		\noindent
		\begin{minipage}[t]{0.50\textwidth}
			% ------------------------------------------------------------------
			% RELATORE
			% ------------------------------------------------------------------
			% Sostituire:
			% - "Chiar.ma Prof." con il titolo corretto del relatore:
			%    * "Chiar.mo Prof." / "Chiar.ma Prof.ssa"
			%    * oppure "Prof." / "Prof.ssa" / "Dott." / "Dott.ssa"
			%   a seconda di quanto concordato con il/la relatore/relatrice.
			% - "NOME RELATORE" con Nome e Cognome del relatore.
			{\large{\bf Relatore:\\
					Chiar.ma Prof.\\
					TODO}} \\[5mm]
			
			% ------------------------------------------------------------------
			% CORRELATORE (FACOLTATIVO)
			% ------------------------------------------------------------------
			% Se HAI un correlatore:
			%   - Sostituisci "Chiar.mo Prof." con il titolo accademico corretto.
			%   - Sostituisci "NOME CORRELATORE" con Nome e Cognome del correlatore.
			%
			% Se NON hai un correlatore:
			%   - Commenta il blocco qui sotto anteponendo % a ciascuna riga,
			%     oppure cancellalo.
			\large{\bf Correlatore: \\
				Chiar.mo Prof. \\
				TODO}
		\end{minipage}
		\hfill
		\begin{minipage}[t]{0.47\textwidth}\raggedleft
			% ------------------------------------------------------------------
			% NOME DELLAUREANDO / LAUREANDA
			% ------------------------------------------------------------------
			% Sostituire "NOME LAUREANDO" con:
			%   Nome Cognome
			% esattamente come compare nei documenti ufficiali.
			{\large{\bf Presentata da:\\
					TODO}}
		\end{minipage}
		
		\vspace{10mm}
		\begin{center}
			% ------------------------------------------------------------------
			% APPELLO, DATA E ANNO ACCADEMICO
			% ------------------------------------------------------------------
			% "NUMERO Appello":
			%   - Sostituire con il numero dell'appello di laurea (es. "I Appello", "II Appello")
			%     secondo le indicazioni della Segreteria/CdS.
			%
			% "Sessione MESE ANNO":
			%   - Sostituire con la data ufficiale della seduta di laurea
			%
			% "Anno Accademico 20xx/20xx":
			%   - Sostituire con l'Anno Accademico corretto (es. "Anno Accademico 2025/2026").
			{\large{\bf Sessione MESE ANNO \\}}
			{\large{\bf Anno Accademico 20TODO/20xx\\}}
		\end{center}
	\end{titlepage}
	\clearpage{\pagestyle{empty}\cleardoublepage}
	\cleardoublepage
	\tableofcontents
	\cleardoublepage
	\listoffigures
	\cleardoublepage
	\listoftables
	\clearpage{\pagestyle{empty}\cleardoublepage}
	\begin{abstract}	
		\setcounter{page}{9}
		\thispagestyle{plain}
		TODO
		IoT (Internet of Things) è l'acronimo utilizzato per indicare la sfera tecnologica che ha alla base l'utilizzo di oggetti intelligenti. 
		Per oggetti intelligenti si intendono quei dispositivi, macchine o attrezzature che si possono connettere a internet per trasmettere, ricevere dati o elaborarli.
		All'interno di questa sfera, uno dei metodi di comunicazione più utilizzati è quello basato su un modello Publish/Subscribe che consente lo scambio di dati tra dispositivi della rete senza la necessità di connessioni dirette tra loro.
		Per funzionare si utilizza il protocollo MQTT, uno dei protocolli più diffusi per lo scambio di messaggi, il quale si basa su una struttura con 3 tipologie di dispositivi connessi: Publisher, Subscriber e Broker.
		I Publisher una volta ottenuti dei dati tramite sensori di vario genere, li pubblica, rendendoli disponibili nella rete. I dati per spostarsi all'interno della rete utilizzano i broker cioè vari punti/server di connessione, i quali permettono il passaggio di dati per poi infine arrivare ai Subscribers per la lettura e l'elaborazione.
		Questa tesi riprende il lavoro effettuato da un collega nell'anno 21-22 approfondendo e sviluppando i casi di test con particolare attenzione a quelle che sono le caratteristiche che 
		possono portare criticità al calcolo della soluzione ottima di percorso per il flusso dei dati.
		La tesi del collega si basa sull'idea di costruire un sistema distribuito sfruttando le proprietà matematiche del rilassamento Lagrangiano, dove ogni nodo conosce solamente i propri vicini.
	\end{abstract}
	
	\chapter{Introduzione}             
	\section{Iot (Internet of Things)}              
	Il termine Iot nasce come acronimo per definire quella sfera tecnologica moderna che si basa sull'utilizzo di dispositivi "intelligenti" che hanno rivoluzionato il modo in cui raccogliamo e processiamo le informazioni sia all'interno di panorami domestici sia nelle industrie.
	Un dispositivo definito intelligente o più sentito comunemente "smart" è un dispositivo in grado di connettersi alla rete per poter svolgere diverse funzioni, tra cui la più importante è lo scambio di informazioni con altri dispositivi.
	Nella mia tesi mi soffermo particolarmente sui "sensori", dispositivi chiave per lo scambio di messaggi in una rete Iot.
	Un sensore è un dispositivo che registra un cambiamento nel mondo fisico, ne genera un messaggio virtuale e poi diffonde il messaggio oppure prende decisioni tramite il valore letto.
	Nella mia tesi vado ad approfondire proprio questo dettaglio.
	Il fatto che un sensore possa generare un messaggio e permettere ad altri dispositivi di prendere decisioni tramite esso significa che se le decisioni devono essere tempestive allora anche il metodo di trasmissione e l'efficienza della stessa lo devono essere.
	\section{Publish and Subscribe}
	Tra le metodologie più famose di trasmissione di messaggi quella su cui mi voglio soffermare è quella del Publish and Subscribe, un paradigma particolarmente adatto all'Iot, nato come sostituzione del classico Client-Server. In questo schema i dispositivi si dividono in Publisher e Subscriber collegati tra loro tramite dei nodi/server di trasmissione definiti Broker.
	In questo modello ognuno a il suo ruolo: I Publisher devono diffondere nella rete le informazioni in loro possesso, i Broker devono ritrasmettere queste informazioni senza permettere danneggiamenti o perdita delle stesse e i Subscriber sono i ricevitori o i richiedenti delle tali.
	Il meccanismo alla base è molto semplice: Il Publisher invia messaggi etichettati con un topic, il broker li riceve e li inoltra solamente ai Subscriber iscritti a quel topic. Questo meccanismo permette un disaccoppiamento tra nodo di partenza e di arrivo in quanto le 2 parti non si conoscono direttamente e non devono per forza essere attivi contemporaneamente.
	Il protocollo standard per la implementazione di questo modello è l'MQTT, progettato per essere leggero e robusto anche su reti instabili.
	
	
	
	
	
	
	
	%Ora vediamo un elenco numerato:         %crea un elenco numerato
	%\begin{enumerate}
	%\item primo oggetto
	%\item secondo oggetto
	%\item terzo oggetto
	%\item quarto oggetto
	%\end{enumerate}
	%
	%\begin{figure}[h]                       %crea l'ambiente figura; [h] sta
	%                                        %   per here, cioè la figura va qui
	%\begin{center}                          %centra nel mezzo della pagina
	%                                        %   la figura
	%%\includegraphics[width=5cm]{figura.eps}%inserisce una figura larga 5cm
	%                                        %se si vuole usare va scommentata
	%%
	%%%%%%%%%%%%%%%%%%%%%%%%%%%%%%%%%%%%%%%%%%inserisce la legenda ed etichetta
	%                                        %   la figura con \label{fig:prima}
	%\caption[legenda elenco figure]{legenda sotto la figura}\label{fig:prima}
	%\end{center}
	%\end{figure}
	%
	%\section{Seconda Sezione}
	%Ora vediamo un elenco puntato:
	%\begin{itemize}                         %crea un elenco puntato
	%\item primo oggetto
	%\item secondo oggetto
	%\end{itemize}
	%
	%\section{Altra Sezione}
	%Vediamo un elenco descrittivo:
	%\begin{description}                     %crea un elenco descrittivo
	%  \item[OGGETTO1] prima descrizione;
	%  \item[OGGETTO2] seconda descrizione;
	%  \item[OGGETTO3] terza descrizione.
	%\end{description}
	%%%%%%%%%%%%%%%%%%%%%%%%%%%%%%%%%%%%%%%%%%crea una sottosezione
	%\subsection{Altra SottoSezione}
	%%%%%%%%%%%%%%%%%%%%%%%%%%%%%%%%%%%%%%%%%%crea una sottosottosezione
	%\subsubsection{SottoSottoSezione}Questa sottosottosezione non viene
	%numerata, ma \`e solo scritta in grassetto.
	%\section{Altra Sezione}                 %crea una sottosezione
	%Vediamo la creazione di una tabella; la tabella \ref{tab:uno}
	%(richiamo il nome della tabella utilizzando la label che ho messo sotto):
	%la facciamo di tre righe e tre colonne, la prima colonna
	%``incolonnata'' a destra (r) e le altre centrate (c):\\
	%\begin{table}[h]                        %ambiente tabella
	%                                        %(serve per avere la legenda)
	%\begin{center}                          %centra nella pagina la tabella
	%\begin{tabular}{r|c|c}                  %tre colonne con righe verticali
	%                                        %   prodotte con |
	%\hline \hline                           %inserisce due righe orizzontali
	%$(1,1)$ & $(1,2)$ & $(1,3)$\\           %& separa le colonne e con
	%\hline                                  %inserisce una riga orizzontale
	%$(2,1)$ & $(2,2)$ & $(2,3)$\\           %  \\ va a capo
	%\hline                                  %inserisce una riga orizzontale
	%$(3,1)$ & $(3,2)$ & $(3,3)$\\
	%\hline \hline                           %inserisce due righe orizzontali
	%\end{tabular}
	%\caption[legenda elenco tabelle]{legenda tabella}\label{tab:uno}
	%\end{center}
	%\end{table}
	%\section{Altra Sezione}\label{sec:prova}%posso mettere le label anche
	%                                        %   alle section
	%\subsection{Listati dei programmi}
	%\subsubsection{Primo Listato}
	%\begin{verbatim}
	%        In questo ambiente     posso scrivere      come voglio,
	%lasciare gli spazi che voglio e non % commentare quando voglio
	%e ci sarà scritto tutto.
	%Quando lo uso è meglio che disattivi il Wrap del WinEdt
	%\end{verbatim}
	%%%%%%%%%%%%%%%%%%%%%%%%%%%%%%%%%%%%%%%%%%non numera l'ultima pagina sinistra
	%\clearpage{\pagestyle{empty}\cleardoublepage}
	%%%%%%%%%%%%%%%%%%%%%%%%%%%%%%%%%%%%%%%%%%per fare le conclusioni
	%\chapter*{Conclusioni}
	%%%%%%%%%%%%%%%%%%%%%%%%%%%%%%%%%%%%%%%%%%imposta l'intestazione di pagina
	%%%%%%%%%%%%%%%%%%%%%%%%%%%%%%%%%%%%%%%%%%aggiunge la voce Conclusioni
	%                                        %   nell'indice
	%\addcontentsline{toc}{chapter}{Conclusioni} Queste sono le
	%conclusioni.\\
	%In queste conclusioni voglio fare un riferimento alla
	%bibliografia: questo \`e il mio riferimento \cite{K3,K4}.
	%%%%%%%%%%%%%%%%%%%%%%%%%%%%%%%%%%%%%%%%%%imposta l'intestazione di pagina
	%\renewcommand{\chaptermark}[1]{\markright{\thechapter \ #1}{}}
	%\appendix                               %imposta le appendici
	%\chapter{Prima Appendice}               %crea l'appendice
	%In questa Appendice non si \`e utilizzato il comando:\\
	%%%%%%%%%%%%%%%%%%%%%%%%%%%%%%%%%%%%%%%%%%\verb"" è equivalente all'
	%                                        %   ambiente verbatim,
	%                                        %   ma si utilizza all'interno
	%                                        %   di un discorso.
	%\verb"\clearpage{\pagestyle{empty}\cleardoublepage}", ed infatti
	%l'ultima pagina 8 ha l'intestazione con il numero di pagina in
	%alto.
	%%%%%%%%%%%%%%%%%%%%%%%%%%%%%%%%%%%%%%%%%%imposta l'intestazione di pagina
	%\chapter{Seconda Appendice}             %crea l'appendice
	%%%%%%%%%%%%%%%%%%%%%%%%%%%%%%%%%%%%%%%%%%imposta l'intestazione di pagina
	%\begin{thebibliography}{90}             %crea l'ambiente bibliografia
	%%%%%%%%%%%%%%%%%%%%%%%%%%%%%%%%%%%%%%%%%%aggiunge la voce Bibliografia
	%                                        %   nell'indice
	%\addcontentsline{toc}{chapter}{Bibliografia}
	%%%%%%%%%%%%%%%%%%%%%%%%%%%%%%%%%%%%%%%%%%provare anche questo comando:
	%%%%%%%%%%%%\addcontentsline{toc}{chapter}{\numberline{}{Bibliografia}}
	%\end{thebibliography}
	%%%%%%%%%%%%%%%%%%%%%%%%%%%%%%%%%%%%%%%%%%non numera l'ultima pagina sinistra
	%\clearpage{\pagestyle{empty}\cleardoublepage}
	%\chapter*{Ringraziamenti}
	%\thispagestyle{empty}
	%Qui possiamo ringraziare il mondo intero!!!!!!!!!!\\
	%Ovviamente solo se uno vuole, non \`e obbligatorio.
\end{document}
